\section{Detection Einheiten der Ladeinfrastruktur}
\label{sec:detection_einheiten_der_ladeinfrastruktur}

Im Folgenden werden die vier Hauptdetektionseinheiten der Testumgebung der Ladeinfrastruktur vorgestellt, die im Rahmen dieses Projekts implementiert wurden und die für den Cyber Threat Report sowie den Aggregierten Cyber Threat Report von zentraler Bedeutung sind. 

    \subsection{System- und Leistungsanomaliedetektion (pi\_anomaly\_detection)}
    \label{subsec:System_und_Leistungsanomaliedetektion}
    Die \textit{Pi Anomaly Detection Einheit} trifft Entscheidungen anhand der folgenden Metriken: Shunt-Spannung, Bus-Spannung, Stromstärke und Leistung. Diese Messdaten werden von den Sensoren der Wallbox erfasst und an Prometheus weitergeleitet. Die beschriebene Anomaliedetektionseinheit greift diese Daten über Prometheus auf und nutzt daraufhin zwei Deep-Learning-Modelle:
    \begin{itemize}
        \item \textbf{xLSTM-Autoencoder:} Dieses Modell operiert auf einem Sliding-Window von $w=44$ Sekunden. Es lernt die Normalverteilung der physikalischen Parameter (Spannung, Strom, Leistung). Eine Anomalie wird detektiert, wenn der Rekonstruktionsfehler $E_t$ den Schwellenwert $\tau = 0.008$ überschreitet. Die Wahrscheinlichkeit wird über eine Sigmoid-Funktion approximiert:
        \[ P_{xlstm} = \frac{1}{1 + e^{-50 \cdot (E_t - \tau)}} \]
        \item \textbf{Transformer-Klassifikator:} Parallel dazu extrahiert das System mittels \textit{tsfresh} 15 statistische Merkmale aus einem 60-Sekunden-Fenster. Ein Transformer-Modell klassifiziert diese aggregierten Merkmale direkt in die Kategorien \textit{benign} oder \textit{attack}.
        \item \textbf{Sequentielles Voting-Verfahren:} Zur Stabilisierung der Detektion nutzt das System einen \textit{Deques}-Puffer der Länge $k=3$. Ein Alarm wird nur ausgelöst, wenn:
        \begin{itemize}
            \item Beide Modelle im aktuellsten Zeitschritt eine Anomalie melden ODER
            \item Mindestens eines der Modelle eine konsistente Sequenz von drei aufeinanderfolgenden positiven Detektionen aufweist.
        \end{itemize}
    \end{itemize}
    
    Die Detektionseinheit reagiert insbesondere auf Stromspitzen und -einbrüche, welche primär über die genannten Metriken identifiziert werden können.

    \subsection{Physikalische Signal- und Temperaturüberwachung (detection\_power\_temperature)}
    \label{subsec:physikalische_signale_und_temperaturueberwachung}

    Diese Detektionseinheit arbeitet mit einer klassischen regelbasierten Logik, welche IDS-Muster zuordnet und auf Basis von Schwellenwerten für physikalische Parameter (Spannung, Leistung, Temperatur) Anomalien detektiert. Die Detektion ist in zwei Bereiche der IDS-Mustererkennung unterteilt:

    \begin{itemize}
        \item \textbf{Temperaturüberwachung:} Diese Anomalieerkennung ruft die Temperaturdaten der letzten 60 Sekunden von Prometheus ab und erkennt Anomalien anhand von drei definierten Regeln:
        \begin{itemize}
            \item \textit{Absolute Obergrenze:} Wenn die Temperatur im abgefragten Zeitfenster einen Wert von 70\,°C überschreitet, wird eine Anomalie detektiert.
            \item \textit{Konstantreihe + Sprung:} Wenn ein Temperatursprung bei einer ansonsten konstanten Reihe eine Toleranz von $\pm 1$ überschreitet, wird eine Anomalie gemeldet.
            \item \textit{Anomalie durch Z-Score:} Wenn die Standardabweichung der Temperaturwerte einen Z-Score von $> 3,0$ erreicht, löst die Anomalieerkennung aus.
        \end{itemize}
        \item \textbf{Wattmeter-Überwachung:} Analog zur System- und Leistungsanomaliedetektion (siehe Abschnitt \ref{subsec:System_und_Leistungsanomaliedetektion}) ruft diese Detection ebenfalls vier Metriken zu Spannung und Leistung der letzten 60 Sekunden ab. Die Detektion prüft, ob diese Werte fest definierte Schwellenwerte (Minimum oder Maximum) unter- bzw. überschreiten. Ist dies der Fall, wird eine Anomalie ausgelöst.
    \end{itemize}

    Sobald die Temperaturüberwachung und/oder die Wattmeterüberwachung eine positive Anomalie detektieren, schlägt auch der kombinierte Alarm \textit{detection\_power\_temperature} aus.

    \subsection{KI-basierte Angriffserkennung im Netzwerk (ids\_power\_ai\_detection)}
    Diese Detection-Komponente analysiert dieselben vier physikalischen Messgrößen, die in Abschnitt~\ref{subsec:System_und_Leistungsanomaliedetektion} beschrieben wurden, verzichtet jedoch vollständig auf statische Schwellenwerte. Stattdessen kommen zwei komplementäre Machine-Learning-Modelle zum Einsatz, die anomales Verhalten im zeitlichen Kontext erkennen:

    \begin{itemize}
        \item \textit{Isolation Forest:} Dies ist ein Verfahren, welches einzelne Messpunkte anhand ihrer statistischen Isolation im Merkmalsraum bewertet. Für jeden Messzeitpunkt wird ein Anomalie-Score berechnet. Weicht ein Datenpunkt signifikant vom gelernten Normalzustand ab, wird dies als Anomalie gewertet.
        \item \textit{LSTM-Autoencoder:} Wie für LSTM-Modelle (Long Short-Term Memory) üblich, betrachtet dieses Modell nicht nur einzelne Messwerte, sondern eine Sequenz von 60 Sekunden (bei einem Abfrageintervall von 5 Sekunden). Das Modell wurde auf normalem Betriebsverhalten trainiert und lernt, typische Zeitreihenmuster zu rekonstruieren. Zur Erkennung wird der Mean Squared Error (MSE) zwischen der originalen Sequenz und ihrer Rekonstruktion berechnet. Überschreitet der MSE des aktuellen Fensters einen empirisch bestimmten Schwellenwert (\texttt{LSTM\_THRESHOLD} = 0,01), wird eine Anomalie ausgelöst.
    \end{itemize}

    \subsection{Hardware-Ereignisüberwachung (detection\_hardwareevents)}
    In diesem Verfahren werden knapp 20 relevante Hardware-Ereignisse kontinuierlich bewertet und mittels eines LSTM-Modells auf anomales Verhalten geprüft. 

    Das Modell analysiert die Daten in einem gleitenden Zeitfenster. Ergibt die Modellvorhersage eine Angriffswahrscheinlichkeit von über 50\,\%, wird das Ereignis als Anomalie klassifiziert. Die Ergebnisse, einschließlich der Einzelwahrscheinlichkeiten für \textit{benign} und \textit{attack}, werden über eine dedizierte Schnittstelle an Prometheus exportiert und zur Verfügung gestellt.
